\section{Experiments}
\label{sec:experiments}

\subsection{Setup}
\label{sec:setup}

\paragraph{Scenes and epochs.}
All main-campaign data comes from BLK360 scans of two industrial indoor
environments that
are, in fact, one room observed across its working life---a property we
exploit deliberately. The \emph{showroom} epoch (June 2026) captures the room
in an exhibition configuration, scanned manually with the door open; its
semantic target is the wall-removed export (38 objects after remnant
filtering). The \emph{robot hall} epoch (July 2026) captures the same room
reorganized as a robot test hall, scanned by autonomous visibility-driven
exploration (81 objects, room-scoped to 55). A third epoch (late July 2026)
adds a controlled-change re-scan at 99.6\% coverage (Sec.~\ref{sec:exp-update}).
Because the two primary epochs differ in layout, furnishing, and scan
modality, they function as independent scenes for per-epoch evaluation
(Secs.~\ref{sec:exp-seg}--\ref{sec:exp-query}) while enabling a real
multi-epoch update study (Sec.~\ref{sec:exp-update}). Scene statistics are
given in Table~\ref{tab:scenes}. A separate scene---a carpeted cafeteria
with sofa booths and vending machines on a different floor of the same
building, scanned in August 2026 (57M points, 28 clusters)---is held out of
the main campaign and used only as a cross-floor, out-of-vocabulary stress
case in Sec.~\ref{sec:exp-ablation}. Both sites are photographed in
Fig.~\ref{fig:photos}.

\begin{figure}[!tp]
\centering
\includegraphics[width=\linewidth]{fig18_scene_photos.jpg}
\caption{The two experimental sites. (a)~The 4F test room, photographed
in its robot-hall configuration (the showroom epoch T1 is an earlier
furnishing of the same room; visible: the AMMR platforms, the
wall-mounted monitor bank of Sec.~\ref{sec:exp-seg}, and the
disinfection robot). (b)~The 8F cafeteria stress scene of
Sec.~\ref{sec:exp-ablation}: fabric sofas, planter partitions, the
vending row, and---center right---the brown door recovered by
escalation on full-resolution renders.}
\label{fig:photos}
\end{figure}

\begin{table}[t]
\centering
\caption{Scene statistics (deterministic detection ``det'' sets, after
structure-remnant filtering).
T1 semantic processing starts from a manually wall-removed export (4.9M
points) of the 46.3M-point raw scan.}
\label{tab:scenes}
\small
\begin{tabular}{@{}lccc@{}}
\toprule
 & Showroom (T1) & Robot hall (T2) & Re-scan (T3) \\
\midrule
Scan mode & manual, 1 setup & exploration & exploration (SOTA), 6 setups \\
Raw points & 46.3M & 68.6M & 70.2M \\
Clean cloud (post-removal) & --- & 267K & 294K \\
Objects (filtered) & 38 & 81 & 79 \\
Room-scoped & 38 & 55 & 41 \\
Removed planes & 10 & 10 & 10 \\
\bottomrule
\end{tabular}
\end{table}

\begin{table}[t]
\centering
\caption{Evaluation-set lineage: how every cluster set used in
Sec.~\ref{sec:experiments} derives from its scan. All counts are
recomputed from the archived stage outputs in the data package;
$\dagger$ marks the one count grounded only in the experiment log.
``conf.''\,=\,GT-confirmed.}
\label{tab:sampleflow}
\footnotesize
\setlength{\tabcolsep}{3.5pt}
\begin{tabular}{@{}ll@{}}
\toprule
Lineage & Stage flow (cluster counts) \\
\midrule
Showroom 57-set (Sec.~\ref{sec:exp-verify}) & 30 pre-split
$\rightarrow$ 57 split $\rightarrow$ 50 after soffit filter
$\rightarrow$ 50 GT-scored (36 conf.) \\
Showroom det set & 45 raw $\rightarrow$ 38 after soffit filter
$\rightarrow$ 38 room-scoped \\
Robot hall det & 88 raw $\rightarrow$ 81 filtered $\rightarrow$ 55
room-scoped $\rightarrow$ 51 GT-scored (32 conf.) \\
T3 det1 & 79 merged $\rightarrow$ 41 room-scoped (= owner audit set)
$\rightarrow$ 26 filter survivors \\
T3 det2 (precut) & 66 $\rightarrow$ 35 room-scoped $\rightarrow$ 27 to
VLM (8 wall sheets excluded) \\
T3 det3/det4 & 33 consolidated $\rightarrow$ 46$^{\dagger}$ adjudicated \\
Cross-epoch KG (present) & rev7$\rightarrow$14:
47, 49, 51, 63, 64, 64, 46, \textbf{46} (current) \\
det5 frozen re-run (Sec.~\ref{sec:exp-endtoend}) & 66 $\rightarrow$ 35
in-room, scored vs.\ the 46-object graph \\
Cafeteria 8F & 28 $\rightarrow$ 25 GT-scored (20 conf.) \\
\bottomrule
\end{tabular}
\end{table}

\paragraph{Ground truth.}
Ground truth was established in two passes: a render-based audit of every
cluster, followed by confirmation from the room's owner. Labels are graded
\emph{confirmed} (owner-verified), \emph{plausible}, or \emph{unknown};
accuracy is reported against confirmed GT (with lenient variants including
plausible where stated), unknown objects are excluded, and structure
artifacts are tracked separately. Two showroom sets must be distinguished:
Table~\ref{tab:scenes} counts the 38-object reproducible det set, while the
verification study (Sec.~\ref{sec:exp-verify}) uses the original 57-cluster
decomposition of the same scan, of which 50 clusters are GT-scored after
excluding 7 structure artifacts---yielding confirmed 36 / plausible 12 /
unknown 2. Table~\ref{tab:sampleflow} connects every evaluation set used
in this section---the 38-object det set, the 50-scored showroom subset,
and the T3 41-, 35-, and 46-object sets---to its derivation stage, so no
denominator has to be inferred from prose. For the room-scoped robot hall (55 objects), GT covers 51
clusters: confirmed 32 (lenient $n{=}42$) after owner corrections. The
cafeteria stress scene follows the same practice: 25 of 28 clusters scored
(3 structure artifacts excluded), confirmed 20 / plausible 5. Both audit
passes share an anchoring channel---the render audit presents proposed
labels that the owner then corrects---so confirmed grades are owner
decisions over physically known rooms rather than blind annotations; the
owner overturned model-favored labels in several cases (a
``refrigerator'' re-graded as a banner, a verified keyboard refuted as a
phantom, ``motors'' re-identified as mobile manipulators), which bounds
but does not eliminate this bias. To measure it directly, a second
annotator---a co-author familiar with the facility who had seen neither
the T3 pipeline output nor the owner audit---labeled all 41 T3 clusters
\emph{blind}, from full-resolution multi-view render sheets carrying only
neutral shuffled IDs, in free vocabulary with a 1--5 confidence. On the
33 clusters both annotators labeled, the blind labels agree with the
owner's on 31 (93.9\%; the two disagreements are low-confidence blind
labels, 2--3/5); the blind annotator abstained on 5 clusters and both
were uncertain on 2. The most informative split for anchoring is by whether the owner had
\emph{overridden} the model: on the 31 clusters where the owner
overturned the pipeline label, the blind annotator independently
reproduced the override on 25 of the 26 it scored (abstaining on 5),
versus 6 of 7 where the owner had accepted the model label. The 25/26 reproduction of the owner's overrides argues against those
corrections being driven primarily by model anchoring; the accepted
subset is small (6/7), so residual anchoring cannot be excluded, and we
therefore describe the owner ground truth as anchoring-bounded rather
than anchoring-free. A third annotation, added for this revision, comes
from an external annotator with no connection to the work and no
familiarity with the facility, who labeled the same 41 sheets blind
under the same protocol. Without site knowledge, agreement with the
owner drops to 21 of the 33 clusters scored (63.6\%; 5 abstentions, 2
both-uncertain), but it is strongly confidence-dependent: 16 of 19
labels given at confidence 4--5 agree (84\%), against 5 of 14 at
confidence $\leq$3, and 11 of the 12 disagreements fall on the two
categories the paper singles out as render-ambiguous --- absorbed or
merged industrial objects (six, e.g.\ the conveyor belt read as a work
table) and ceiling or wall remnants read as floor-level objects (five).
We read the external
result as the honest bound on render-only annotation: distinctive
furniture and robots are labeled consistently by anyone, while the
categories that required on-site knowledge in the owner audit are
exactly those an unfamiliar annotator cannot resolve from renders. The
blind kits, all three answer sheets,
and scoring are in the data package.

\paragraph{Metrics.}
Classification accuracy is top-1 against GT under a documented
synonym-tolerant matching (e.g., \emph{mobile manipulator} $\approx$
\emph{robotic arm}). Verification quality is reported as per-tier precision
(unanimous, escalated, unverified). For query evaluation, a
\emph{hallucination} is an unsupported atomic factual claim, and the rate is
unsupported claims over all factual claims in the answer set. Update quality
uses the metrics of Sec.~\ref{sec:matching}: matching accuracy,
identity-switch rate, false-new, false-absent, and unchanged stability.

\paragraph{Controls.}
The verifier for the main campaign (Secs.~\ref{sec:exp-seg}--\ref{sec:exp-update}) is fixed to a single
model, \texttt{claude-\hspace{0pt}sonnet-\hspace{0pt}4-6} (the frontier
mid-tier at campaign start):
every verification, escalation, and owner interaction in the accumulated
knowledge-graph lineage was produced under this one verifier, so results
remain comparable across sections and revisions. The verifier choice itself
is then evaluated in Sec.~\ref{sec:exp-ablation}, which repeats the
identical Stage-B protocol across five commercial VLMs and re-runs the
headline verification study under the best-ranked one.
All query conditions share the same LLM (\texttt{claude-sonnet-4-6}), system
prompt, query set, decoding parameters, retry policy, and forced answer
schema; conditions differ \emph{only} in the knowledge view they receive.
The LLM-only condition---an \emph{unverified Stage-A LLM} baseline---receives the same scene's Stage-A object dump, so it
is not information-starved---it is verification-starved.
Threshold provenance is stated in two parts. Constants in the base
pipeline (Secs.~\ref{sec:preprocess}--\ref{sec:upsert}) were fixed on
T1/T2 and applied unchanged to T3 and the cafeteria. The corrective-layer
thresholds of Sec.~\ref{sec:corrective} were developed on the test-room
lineage---including failures revealed by the T3 owner audit
(Sec.~\ref{sec:exp-owner})---and were then frozen before evaluation on
the cafeteria stress scene; accordingly, T3 is treated as a retrospective
development-scene replay of those layers rather than a held-out
evaluation (Sec.~\ref{sec:exp-endtoend}). Nothing is trained, so the
split that matters is the development room (T1/T2/T3) versus the
held-out cafeteria, and Sec.~\ref{sec:exp-endtoend} states which claims
survive on which side of it.

\paragraph{Generative AI use.}
Generative models enter this study in two strictly separated roles,
disclosed here per the journal's GenAI policy. As \emph{experimental
components under test}: commercial VLM/LLM services perform Stage-B
multi-view verification (\texttt{claude-sonnet-4-6} throughout the main
campaign; Anthropic PBC, San Francisco, CA, USA), the query-benchmark
answering under all four conditions (the same model), place-name
generation from ring codes, and the five-verifier comparison of
Sec.~\ref{sec:exp-ablation} (Anthropic and OpenAI models); every such
call is made under a forced tool-use schema with prompts, decoding
parameters, model identifiers, and per-call usage archived in the data
package, and these models contributed no text to this article. As a
\emph{writing aid}: a generative AI assistant was used for language
editing and LaTeX formatting of the manuscript, as declared in the
Acknowledgments; it played no role in study design, data collection, or
analysis beyond the archived experimental calls above.

\subsection{Segmentation: Repeatability, Filtering, and Sensitivity}
\label{sec:exp-seg}

\paragraph{Repeatability.}
Running the full pipeline twice from the same E57---including ten-plane
structure removal---reproduces every cluster byte-identically on both
scenes (hash-verified). To probe robustness beyond one software stack, we
packaged the backbone into a version-pinned container and re-ran the
robot-hall scan under a deliberately different environment (glibc 2.41
vs.\ 2.39, Python 3.12.14 vs.\ 3.12.3): all 81 retained clusters remain
byte-identical to the native run, and consecutive container runs
reproduce each other exactly. Running the same container on a second
machine with a different CPU vendor and microarchitecture (Intel Core
i9-13900H laptop vs.\ the AMD Ryzen~9 9900X reference workstation)
reproduces the identical per-cluster checksums and aggregate hash---the
determinism claim therefore holds across operating environments and
x86-64 CPU vendors under the pinned library set. The verification kit
(container recipe,
reference checksums, comparison script) ships with the data package so
the check can be repeated elsewhere; bit-level
reproducibility across instruction-set architectures is not claimed.
This is the property the incremental machinery
builds on: identical clusters mean carryable verification results and
attributable diffs.

\paragraph{Remnant filter.}
The structure-remnant filter of Sec.~\ref{sec:decomposition} removed all
fourteen owner-confirmed ceiling-soffit phantoms (seven per scene: showroom
45$\rightarrow$38, robot hall 88$\rightarrow$81) while leaving every retained
cluster
unchanged. Before the filter, these phantoms did not merely inflate the
object count: they produced confidently verified fictions (keyboard panels,
handrails) that survived multi-view voting, since every view showed the same
wall-colored slab. Detecting structure remnants \emph{geometrically},
upstream of semantics, proved more reliable than asking the verifier to
doubt itself (cf. the pose-context ablation, Sec.~\ref{sec:exp-ablation}).

\paragraph{Parameter sensitivity.}
A 36-combination sweep (giant-split footprint $\times$ DBSCAN
$\varepsilon$ $\times$ minimum points) shows a stable plateau
at footprint 2.0--2.5\,m: instance recall degrades at 3.0\,m as neighbors
merge and at 1.5\,m as equipment shatters, while
$\varepsilon$/min-point extremes collapse recall to 0.21. Because scoring
a sweep against the operating point's own decomposition would be
circular, we re-scored every combination against an independent
object-level reference: the owner-audited inventory, in a
membership-independent form (all clean-cloud points falling inside each
owner-audited oriented box, so no clustering decision defines the
reference). On the robot hall (42 reference objects) the conclusion
survives intact: the min-points-100 band is a broad plateau, and the
operating point (2.5\,m / 0.30 / 100) sits at the grid maximum under
the independent reference (recall 0.952, a uniform $\sim$0.05 below the
cluster-referenced numbers, with the plateau shape unchanged). On the
showroom, the audit surfaced an error we report plainly: the original
showroom sweep rows had joined the reference onto the operating run's
own clusters, and against the genuinely audited 57-object inventory
\emph{no} grid combination of the base clustering exceeds recall 0.33
--- the audited granularity of that scene is produced by the
split-refinement stage, not by any base-DBSCAN setting. The showroom
rows are therefore evidence for the necessity of the refinement stage,
and the parameter-robustness claim rests on the non-circular robot-hall
table (full grids in the data package).

\paragraph{Conditional Point-SAM refinement.}
Under a pre-specified rule (adopt iff $\geq$50\% of merged-cluster failures
separate correctly \emph{and} label accuracy improves $\geq$5 points),
Point-SAM separated 3/6 flagged merged clusters and sub-cluster labels
reached 9/11 versus 5/11 for their parents (+36.4 points,
Fig.~\ref{fig:baselines}), so the rule is met for \emph{flagged merged
candidates}; because automatic candidate selection also produced 2/6
false positives that fragment single objects, Point-SAM is retained only
as an optional, manually triggered diagnostic refinement---it is not part
of the frozen pipeline of Sec.~\ref{sec:corrective} and contributes
nothing to the end-to-end replay of Sec.~\ref{sec:exp-endtoend}.

\begin{figure}[t]
\centering
\includegraphics[width=0.95\linewidth]{fig13_pointsam_uni3d.png}
\caption{Learned components used as bounded tools. (a)~The Uni3D provisional
label degrades as the vocabulary narrows to the deployment's true classes
(Sec.~\ref{sec:exp-ablation})---why it is treated as untrusted context.
(b,c)~Conditional Point-SAM refinement of a flagged merged cluster: a
machine, an office chair, and a chair re-emerge from a single
\texttt{clutter} blob, each re-verified independently.}
\label{fig:baselines}
\end{figure}

\paragraph{Coverage stress case.}
The T3 epoch supplied a stress case: at 99.6\% coverage, inter-object
gaps close and DBSCAN chains 120K points (41\% of the clean cloud) into one
$10.4\times9.6$\,m mega-cluster that survives the giant-split. Refinement
recovers little of the resulting diff inflation---of the epoch's 45 absent
nodes, 43 retain a source cluster to re-associate against; Point-SAM
($K{=}7$ prompts) re-associates 4 of these 43, tighter re-splitting
($\varepsilon{=}0.12$) 5 of 43---which is itself informative: $\sim$88\% of
that epoch's absences reflect genuine scene change, not segmentation
artifacts (Sec.~\ref{sec:exp-update}).

\paragraph{Backbone selection.}
The DBSCAN\,+\,Uni3D backbone was chosen against learned alternatives on
a wall-removed industrial scene (122.7K points), summarized in
Table~\ref{tab:backbones}. The metrics are deliberate proxies (no point-level
ground truth exists for this scene): point coverage, the fraction of
objects the open-vocabulary classifier cannot name (a boundary-quality
proxy), and label diversity. ScanNet-pretrained SPFormer hallucinates
furniture on industrial content (67 of 88 instances ``chair'') at half
the coverage; PTv3's closed-set semantic head labels 46\% of a scene
\emph{whose walls were already removed} as \texttt{wall} --- a closed
set cannot abstain; and PointCLIP~V2 collapses to a single class.
The geometric backbone wins not by being smarter but by refusing to
guess: its failure mode is unlabeled coverage, which the verification
and owner layers are built to absorb. Point-SAM enters only as the
conditional refiner above.

\begin{table}[t]
\centering
\caption{Segmentation$\times$classification backbone comparison on a
wall-removed industrial TLS scene. The metrics are screening proxies
(no point-level GT) and do not measure instance quality; the PTv3 row
is a semantic-segmentation head whose outputs are not
instance-comparable with the clustering rows and is included only to
document why closed-set semantic heads were screened out. Instance
quality of the retained backbone is validated independently in
Secs.~\ref{sec:exp-seg} (owner-referenced sweep)
and~\ref{sec:exp-endtoend} (end-to-end recall).}
\label{tab:backbones}
\footnotesize
\setlength{\tabcolsep}{4pt}
\begin{tabular}{@{}lcccl@{}}
\toprule
Pipeline & Coverage & Unnameable & Classes & Dominant failure \\
\midrule
\textbf{DBSCAN + Uni3D} & \textbf{98.3\%} & \textbf{23.3\%} & 13 &
unlabeled coverage \\
SPFormer + Uni3D & 51.1\% & 45.9\% & 9 & furniture bias (chair 67/88) \\
PTv3 (semantic) & --- & --- & 17/20 emitted & 46\% ``wall'' after removal \\
DBSCAN + PointCLIP~V2 & 98.3\% & --- & 1--7 & collapse (pipe 23/30) \\
\bottomrule
\end{tabular}
\end{table}

\paragraph{Diffuseness-gated re-splitting.}
This layer and the next are defined, with triggers and parameters, in
Sec.~\ref{sec:corrective}; here we report what they did.
Two T3 clusters that survived every upstream filter are \emph{diffuse
webs}: 95K and 13K points spread over 80 and 12\,m$^2$ at a 2D occupancy
fill below 0.3, versus 0.5--0.85 for every real object in the scene ---
a clean geometric signature, but one density clustering cannot act on,
because the webs are connected at $\sim$1.8\,cm point spacing at any
usable $\varepsilon$. A two-stage re-split resolves them by \emph{cause}
instead of size: points hugging the exploration-grid wall line (inside
the \emph{room polygon}, so wall-flush furniture survives the test that
free-space masks would fail) are structure noise ($\sim$70\% of the
$\sim$95K-point residual); out-of-room bleed-through accounts for another
7\%; the remainder is re-clustered by occupancy-core erosion and, for
compounds, split at the support plane detected as the $z$-histogram
peak 0.5--0.95\,m above floor. The owner's ground truth for the largest
compound --- ``a robot arm on a desk, part of a mobile robot, wall
noise'' --- is recovered exactly: the arm (consensus-verified), the
desk (slab plus under-desk storage), and the partially scanned mobile
robot, plus three columns and a plant. Applied to a $6.2$\,m
wall-flush cluster that verification had confidently labeled
\texttt{electrical cabinet} (0.83), the same machinery --- wall-band
strip, support plane, then density-valley cuts along the wall in the
panel band --- recovers the owner's inventory at unit granularity:
nine of the ten owner-confirmed monitor units, the desk beneath them,
and, from the off-wall remainder, a wheelchair robot and a single
table leg. Unit boundaries are approximate --- the dark screens return
few points, so the cuts follow density valleys in the surrounding wall
band --- and the result is claimed as unit-level recall, not precise
per-unit segmentation. The gate turns ``too diffuse to be an object''
from a discard decision into a decomposition trigger.

\paragraph{Structure-condition ensemble.}
Structure removal is itself a segmentation parameter: it amputates
wall-flush objects along with the wall. Running the geometric pass
under three removal conditions (none / ceiling only / full) and fusing
the results recovered two multi-metre tables that exist in \emph{no}
single-condition segmentation of the production path --- they appear
only when walls are left in place --- raising in-room recall by four
real objects ($\sim$8\%) as a standalone pass (its incremental gain in
the cumulative stack of Table~\ref{tab:layers} is smaller because the
re-split and wall-compound layers already recover part of the same set). The fusion also quantifies label fragility:
only 10 of 37 re-found objects keep the same label across conditions,
and cross-condition \emph{voting} corrected nothing, because the
verifier's biases are correlated across conditions (both recovered
tables were unanimously mis-verified as conveyor belts in every
condition, while the room's actual conveyor unit stayed
\texttt{clutter} in every condition). The ensemble is therefore used as
a recall tool plus a per-object \texttt{labelStability} attribute that
prioritizes owner review --- not as a label fuser.

\subsection{Verification: Fusion and Escalation}
\label{sec:exp-verify}

\begin{table}[t]
\centering
\caption{Verification accuracy against confirmed GT, and per-tier precision.
Accuracy rows are over confirmed GT ($n{=}36$/$32$); trigger rates and tier
counts are over all GT-scored clusters (50 showroom / 51 robot hall), so,
e.g., 21 unanimous $+$ 29 split $=$ 50.}
\label{tab:verification}
\small
\begin{tabular}{@{}lcccc@{}}
\toprule
 & \multicolumn{2}{c}{Showroom ($n{=}36$)} & \multicolumn{2}{c}{Robot hall ($n{=}32$)} \\
 & base-4 & escalated & base-4 & escalated \\
\midrule
Accuracy & 77.8\% & \textbf{88.9\%} & 78.1\% & \textbf{84.4\%} \\
95\% Wilson CI & 61.9--88.3 & 74.7--95.6 & 61.2--89.0 & 68.2--93.1 \\
Recovered / regressed & \multicolumn{2}{c}{5 / 1} & \multicolumn{2}{c}{3 / 0} \\
Escalation trigger & \multicolumn{2}{c}{58\% (29/50)} & \multicolumn{2}{c}{51\% (26/51)} \\
\midrule
Unanimous precision & \multicolumn{2}{c}{100\% (19/19; 21 unan.)} & \multicolumn{2}{c}{94.4\% (17/18; 25 unan.)} \\
Escalated precision & \multicolumn{2}{c}{100\%} & \multicolumn{2}{c}{88.9\%} \\
Unverified precision & \multicolumn{2}{c}{60\%} & \multicolumn{2}{c}{40\%} \\
\bottomrule
\end{tabular}
\end{table}

Escalation lifts confirmed-GT accuracy from 77.8\% to 88.9\% on the showroom
and from 78.1\% to 84.4\% on the robot hall (lenient: 76.2\%
$\rightarrow$ 83.3\%, $n{=}42$), recovering 5 and 3 split-vote objects while
regressing 1 and 0---a consistent observed improvement across both
configurations, though at these sample sizes the paired contrast does not
reach conventional significance (exact McNemar on recovered/regressed:
$p{=}0.219$ and $p{=}0.25$), which the Wilson intervals of
Table~\ref{tab:verification} make explicit (Table~\ref{tab:verification};
Fig.~\ref{fig:esccost} shows the accuracy--cost trade-off at
$\sim$\$0.009/call; escalation adds 232 and 208 calls per scene). The tier
structure
justifies gating: unanimous and escalated tiers reach 88.9--100\% precision,
while the unverified residue is 40--60\%---exactly the content the query
interface must demote. These precision figures are scoped to confirmed
GT: on the showroom, one unanimous vote falls on a \emph{plausible}-graded
cluster and is wrong there (a suspected mobile robot absorbed as
\texttt{clutter}), so the tier's precision over confirmed$+$plausible GT is
20/21 rather than 19/19---the scoping convention is reported so that the
headline is not read as a general property of unanimity (the stress scene
of Sec.~\ref{sec:exp-ablation} degrades it substantially). A repeat of the showroom run under identical
settings bounds the verifier's run-to-run variance: base-4 accuracy
reproduces exactly (77.8\%), escalated accuracy shifts by a single object
(86.1\% vs.\ 88.9\%), the trigger rate by two points (60\% vs.\ 58\%), and
unanimous and escalated precision hold at 100\% in both runs---a
$\pm$3-point scale against which the verifier comparisons of
Sec.~\ref{sec:exp-ablation} should be read. A separate check evaluates the room-scoped T1 objects
in the update-study frame (Sec.~\ref{sec:exp-update}), where showroom GT is
transferred to the re-decomposed clusters by bidirectional point overlap;
on the resulting small confirmed set ($n{=}8$), verified-tier precision is
100\% with 7/7 unanimity.

Verification yield tracks scan quality: the manual scan leaves 47\% of
objects unverified versus 24\% for exploration scans, and re-processing the
T3 scan from a single setup (1/6 coverage) drops the verified tier from
75.6\% to 56.6\% on identical geometry. Coverage buys verifiability.

\paragraph{Per-tier recall.}
Precision alone understates what gating costs, so
Table~\ref{tab:tierrf} completes the tier picture with recall and $F_1$
against confirmed GT. The pattern is consistent: the combined verified
tier reaches recall 0.722/0.781 on the two industrial scenes
($F_1$ 0.839/0.847) but only 0.250 on the cafeteria, where the gate
abstained on half of the confirmed objects---the stress scene's
segmentation failures surface here as recall, not as precision. The
escalated tier is small everywhere (recall 0.194--0.250), consistent
with its role as a recovery channel rather than a primary path.

\begin{table}[t]
\centering
\caption{Per-tier precision, recall, and $F_1$ on confirmed ground truth
(synonym-tolerant). Recall counts a confirmed object as retrieved only if
the tier contains it with the correct type; the unverified row is the
abstention tier, so its ``precision'' is only the fraction of abstained
objects whose withheld top vote happened to be correct.}
\label{tab:tierrf}
\footnotesize
\setlength{\tabcolsep}{3.5pt}
\begin{tabular}{@{}llrrrrr@{}}
\toprule
Scene & Tier & $n$ & $n_{\mathrm{conf}}$ & Prec. & Rec. & $F_1$ \\
\midrule
Showroom & verified (unanimous) & 21 & 19 & 1.000 & 0.528 & 0.691 \\
 & verified\_escalated & 8 & 7 & 1.000 & 0.194 & 0.326 \\
 & combined verified & 29 & 26 & 1.000 & 0.722 & 0.839 \\
 & unverified (abstained) & 21 & 10 & 0.600 & 0.167 & 0.261 \\
 & \emph{confirmed GT total} & 50 & 36 & & & \\
\midrule
Robot hall & verified (unanimous) & 25 & 18 & 0.944 & 0.531 & 0.680 \\
 & verified\_escalated & 16 & 9 & 0.889 & 0.250 & 0.390 \\
 & combined verified & 41 & 27 & 0.926 & 0.781 & 0.847 \\
 & unverified (abstained) & 10 & 5 & 0.400 & 0.062 & 0.108 \\
 & \emph{confirmed GT total} & 51 & 32 & & & \\
\midrule
Cafeteria & verified (unanimous) & 8 & 7 & 0.571 & 0.200 & 0.296 \\
 & verified\_escalated & 6 & 3 & 0.333 & 0.050 & 0.087 \\
 & combined verified & 14 & 10 & 0.500 & 0.250 & 0.333 \\
 & unverified (abstained) & 11 & 10 & 0.200 & 0.100 & 0.133 \\
 & \emph{confirmed GT total} & 25 & 20 & & & \\
\bottomrule
\end{tabular}
\end{table}

\paragraph{Threshold sensitivity.}
Because the 0.6 share threshold gates status only
(Sec.~\ref{sec:verification}), its effect can be quantified exactly by
re-scoring the stored votes at alternative thresholds---no new verifier
calls. Sweeping 0.40--0.90 in steps of 0.05 shows 0.6 sitting on a
plateau rather than a cliff. On the showroom, thresholds 0.55 and 0.60
produce identical tiers (escalated-tier precision 100\%); lowering to
0.50 trades combined verified-tier precision (1.000\,$\rightarrow$\,0.966)
for recall (0.722\,$\rightarrow$\,0.778), and raising to 0.70 keeps
precision saturated while dropping recall to 0.639. The robot hall
behaves the same way (at 0.60: escalated precision 0.889, recall 0.781;
at 0.70: precision 1.000, recall 0.688), and on both industrial scenes
F1 is maximized on a broad 0.45--0.65 band with every 0.05 step moving
at most one or two objects. The stress scene shifts the whole curve
downward without changing its shape---the threshold cannot repair
render- or segmentation-level failures
(Sec.~\ref{sec:exp-ablation})---which is the scene dependence noted for
the tiers themselves. Between 0.50 and 0.70 only 11/34, 12/49, and 7/19
escalated objects change tier on the three scenes, and those bands mix
correct and incorrect labels almost evenly, which is why no alternative
threshold dominates 0.6 by more than single-object noise. We therefore
retain the a-priori 0.6 rather than tuning it post hoc.

\paragraph{Error taxonomy.}
Four recurring error classes emerged (Fig.~\ref{fig:taxonomy}):
(a) \emph{structure-remnant phantoms}---soffit bands verified as keyboards;
eliminated upstream by the remnant filter;
(b) \emph{unanimous-but-wrong}---a hanging chain unanimously verified as a
door handle; an escalated keyboard the owner refuted; partition fragments as
chairs; at T3, a residual $7.4$\,m wall segment that survived the remnant
filter was escalation-verified as a staircase---owner-refuted (``just the
wall'')---and its phantom label then propagated into a place name
(\texttt{staircase\_landing}, Sec.~\ref{sec:place}), showing that place
naming is verifier-bounded too and that wall-scale remnants can slip past a
filter tuned for ceiling-soffit bands (cf. Sec.~\ref{sec:discussion}).
The same propagation reaches the relation layer: a keyboard verified at
full confidence in a room that contains no keyboard (owner-refuted) had
been auto-designated its place's key object before the refutation demoted
it---key-object designation therefore honors owner feedback
(Sec.~\ref{sec:upsert}). A third refutation is the mirror image: the
room's large wall-mounted TV was geometrically detected
(\texttt{monitor\_006}, $2.5\times1.8$\,m) yet left \texttt{unverified}
at 0.26 confidence, while a small co-located cluster was
escalation-verified as a fire extinguisher (0.72) the owner refuted---the
salient real object failed the gate a phantom passed.
Unanimity is strong evidence, not proof, and the gate inherits the
verifier's blind spots;
(c) \emph{clutter absorption}---a disassembled robot, a poster stand, and a
low-reflectance floor fan (674 points) each unanimously verified as
\texttt{clutter}: geometrically detected, semantically flattened;
(d) \emph{render ambiguity} $\rightarrow$ \texttt{unverified}, the intended
failure mode.
Classes (b) and (c) motivate reporting verification as \emph{tiered
evidence} rather than truth, and drive the structural gating design.

\paragraph{Class-level breakdown.}
Table~\ref{tab:classdist} disaggregates every scene's ground truth into
class families and per-family verification outcomes, making the error
classes above quantitative. Three patterns stand out. First, the
vocabulary gap is class-structured: robots are the worst family
everywhere they occur (0/4, 0/3, and 0/8 correct across the three
robot-bearing sets)---absorption, not misperception, since most land in
\texttt{clutter} or abstain. Second, displays and fixtures abstain
heavily on the showroom (0/5 and 0/6 correct, all withheld or
mislabeled)---flat panels and ceiling fixtures are exactly the
render-ambiguous shapes the taxonomy predicts. Third, the T3 owner audit
shows the same distribution end-to-end (6 correct / 25 wrong / 10
abstained over 41), confirming that per-family results, not the
aggregate, are the actionable signal for vocabulary and geometry work.

\begin{table}[t]
\centering
\caption{Ground-truth object classes per scene, grouped into families,
with each object's verification outcome in the canonical run
(synonym-tolerant matching as in the escalation analysis). ``Filtered''
= excluded before scoring (structure artifact or out of room scope).
For the T3 owner audit, correctness follows the owner's category
labels; one verified cluster with no owner answer is counted as wrong
in the total row only.}
\label{tab:classdist}
\footnotesize
\setlength{\tabcolsep}{3.5pt}
\begin{tabular}{@{}llrrrrr@{}}
\toprule
Scene & Family & GT & Corr. & Wrong & Unver. & Filt. \\
\midrule
Showroom (57-set) & seating & 5 & 4 & 1 & 0 & 0 \\
 & robots & 4 & 0 & 1 & 3 & 0 \\
 & displays/monitors & 5 & 0 & 0 & 5 & 0 \\
 & machines/electrical & 6 & 3 & 0 & 3 & 0 \\
 & furniture & 2 & 0 & 0 & 2 & 0 \\
 & fixtures & 6 & 0 & 0 & 6 & 0 \\
 & structure-noise & 7 & 0 & 0 & 0 & 7 \\
 & clutter/unknown & 22 & 20 & 0 & 2 & 0 \\
 & \emph{total} & 57 & 27 & 2 & 21 & 7 \\
\midrule
Robot hall (det) & seating & 5 & 4 & 0 & 1 & 0 \\
 & robots & 3 & 0 & 1 & 2 & 0 \\
 & displays/monitors & 1 & 0 & 0 & 0 & 1 \\
 & machines/electrical & 1 & 0 & 0 & 0 & 1 \\
 & furniture & 2 & 0 & 0 & 0 & 2 \\
 & fixtures & 8 & 3 & 1 & 4 & 0 \\
 & structure-noise & 8 & 0 & 0 & 1 & 7 \\
 & clutter/unknown & 53 & 25 & 7 & 2 & 19 \\
 & \emph{total} & 81 & 32 & 9 & 10 & 30 \\
\midrule
Cafeteria & seating & 10 & 1 & 4 & 5 & 0 \\
 & machines/electrical & 1 & 1 & 0 & 0 & 0 \\
 & furniture & 1 & 0 & 1 & 0 & 0 \\
 & fixtures & 5 & 0 & 1 & 4 & 0 \\
 & structure-noise & 3 & 0 & 0 & 0 & 3 \\
 & clutter/unknown & 8 & 4 & 2 & 2 & 0 \\
 & \emph{total} & 28 & 6 & 8 & 11 & 3 \\
\midrule
T3 re-scan (owner GT) & seating & 5 & 2 & 2 & 1 & 0 \\
 & robots & 8 & 0 & 6 & 2 & 0 \\
 & displays/monitors & 2 & 0 & 1 & 1 & 0 \\
 & machines/electrical & 4 & 0 & 4 & 0 & 0 \\
 & furniture & 4 & 0 & 3 & 1 & 0 \\
 & fixtures & 4 & 0 & 2 & 2 & 0 \\
 & structure-noise & 9 & 0 & 6 & 3 & 0 \\
 & clutter/unknown & 5 & 4 & 0 & 0 & 0 \\
 & \emph{total} & 41 & 6 & 25 & 10 & 0 \\
\bottomrule
\end{tabular}
\end{table}

\subsection{Owner Audit and Corrective Feedback (T3)}
\label{sec:exp-owner}

The preceding subsection measured what the pipeline believes; this one
measures what the room's occupant knows, and how that knowledge feeds
back into the graph. It exercises the revision and upsert machinery ahead
of its formal evaluation (Sec.~\ref{sec:exp-update}); only the feedback
mechanics matter here.

\paragraph{Exhaustive audit and mitigations.}
The room occupant exhaustively labeled all 41 T3 clusters, giving a
distinct T3 owner-audited ground-truth set beyond the showroom scoring
sets.
Of the 31 verified-tier clusters, 6 were correct, 13 were real objects
flattened to \texttt{clutter}, 6 were structure-noise phantoms, 2 carried
a wrong specific label, and 3 were multi-object merges (one unanswered).
The absorption cause is now identifiable: the room is a robotics testbed
(mobile robots, manipulators, charging stations, a disinfection robot),
and none of these classes exist in the classifier vocabulary. The
unverified tier behaved as designed in five structure/fixture cases but
also withheld three real objects, including the wall TV of
Sec.~\ref{sec:exp-verify}. Three audit-driven mitigations follow; their
definitions, triggers, and parameters are given in
Sec.~\ref{sec:corrective} (layers 4--5), and this subsection reports their
development-time evaluation.
\emph{(i) Hybrid structure filter:} registering every cluster against the
exploration-scan occupancy grid --- an independent, drift-free wall prior
--- and combining a wall-line hug test with local (per-metre-segment)
sheet thickness, a full-height band rule ($z$-span $>2.5$\,m), a ceiling
band, and a room-polygon outside test removes structure noise at 12/12
recall with 3 false positives (thin wall-flush real objects: a cabinet
leg, a poster stand). Crucially, the wall-flush monitor bank survives:
its wall-hug fraction is 99\% (the grid map draws it as wall), but its
1.8\,m height and per-segment depth separate it from the 3.2\,m
floor-to-ceiling wall band that plane removal missed. All seven verified
structure phantoms --- the six structure-noise cases above plus a
fluorescent fixture verified as a keyboard at 3.3\,m (audited as a wrong
label) --- are eliminated upstream.
\emph{(ii) Recovery re-verification:} re-asking the verifier with
\texttt{clutter} treated as \emph{unidentified} plus a candidate
shortlist recovered 4 specific types at 4/5 precision with 12 honest
keeps; injecting site vocabulary (robot classes) as candidates doubled
recovery (9 correct) but halved precision (8 wrong) through candidate
anchoring --- the TV became a ``conveyor belt'' when the shortlist ranked
it first. \emph{(iii) Agreement gating:} adopting a recovered label only
when both passes agree, or when a pass agrees with the cluster's own
base label, yields five adoptions at 5/5 precision (two chairs, an
office chair, a machine, and the TV) and resurrects no owner-refuted
phantom. Domain knowledge helps recall; only agreement makes it safe.
Fig.~\ref{fig:structfilter} shows the filter's effect on the T3 epoch.
Regression on the other epochs (no ground truth, manual inspection)
flags 8/50 T1 and 18/55 T2 clusters --- dominated by the documented
soffit phantoms (keyboards, chairs at 2.4\,m, a stair at 3.4\,m) and a
cluster the classifier itself had named \texttt{wall} --- while every
door, chair, machine, elevator, and wall-mounted control panel survives;
the one systematic casualty is a closed-door fragment lying in the wall
plane, consistent with the door-detection asymmetry of
Sec.~\ref{sec:discussion}. Finally, refutations, corrections, and
adopted labels are applied as a fourth KG revision through the same
upsert history mechanism as a rescan: eight phantoms leave the graph,
the poster stand is corrected in place, five recovered types land, the
TV becomes its region's key object, and the relation layer sheds nine
noise-anchored edges (25 to 16 intra-place). Re-deriving ring codes and
names from the revised graph heals the place layer end-to-end: the
phantom-derived \texttt{staircase\_landing} becomes
\texttt{display\_briefing\_area} (TV + poster stand), the
keyboard-derived \texttt{workstation\_area} becomes
\texttt{machine\_area}, and the seating region's recovered chairs turn
\texttt{cluttered\_seating\_area} into \texttt{seating\_area} ---
owner feedback propagates through all three layers of the model: object,
relation, and place.

\paragraph{Feedback at scale: bidirectional, corrective, and
time-variant.}
Extending the audit loop through nine further KG revisions shows owner
feedback operating in all four directions the model supports.
\emph{Restoration}: an object the owner had refuted as absent was later
declared present again (the room changed between audits); a
color-seeded carve at the owner-given corner recovered a 319-point red
cylinder that two-pass verification typed \texttt{fire extinguisher}
(0.72) --- the original failure was segmentation (never proposed), not
classification, and owner ground truth is itself \emph{time-variant},
which is why feedback is recorded as highest-trust evidence with
history rather than as immutable truth. \emph{Refutation}: a
raw-condition ensemble find was refuted as wall paint --- the patches
surround the room's wall-mounted TV, whose glossy black screen returns
almost no lidar, so the ``object'' was the halo around a hole.
\emph{Label correction}: the correlated-bias tables of
Sec.~\ref{sec:exp-seg} and a poster stand (tie-break ``fire
extinguisher'', 0.72) were corrected in place, each with the VLM's
verdict preserved in provenance. \emph{Attribute
correction}: wall-remnant clutter was marked immovable in the implicit
layer. Adjudication itself gained three inexpensive refinements along
the way: canonical synonym matching (``robot\_arm'' vs ``robot
manipulator'' is agreement, not disagreement), a focused tie-break pass
that re-asks with exactly the two disputed labels (resolving desk vs
machine at 0.62), and a coarse three-level height attribute
(low/mid/high relative to the 2D-lidar scan plane and ceiling) injected
as a hard physical constraint --- the targeted successor to the
pose-context defense of Sec.~\ref{sec:exp-ablation}, and now a
queryable implicit attribute of every node.

\begin{figure}[t]
\centering
\includegraphics[width=\linewidth]{fig14_structure_filter.png}
\caption{Hybrid structure-noise filter on the T3 epoch. (a) All 41
detected clusters; the 15 flagged ones are colored by rule --- wall
remnant/band (grid-map hug $\times$ thin vertical sheet, or long
full-height band), ceiling fixture/noise, outside the room boundary.
(b) The 26 survivors with recovery / owner-corrected labels. The
wall-flush monitor bank (labeled \texttt{tv}) and wall-mounted equipment
survive; residual merge labels (\texttt{electrical cabinet},
\texttt{ducting}) are retained pipeline output, not corrections.}
\label{fig:structfilter}
\end{figure}

\begin{figure}[t]
\centering
\includegraphics[width=0.85\linewidth]{fig6_escalation_cost.png}
\caption{Escalation accuracy--cost curves for both scenes.}
\label{fig:esccost}
\end{figure}

\begin{figure}[t]
\centering
\includegraphics[width=0.9\linewidth]{fig10_error_taxonomy.png}
\caption{Verification error taxonomy: remnant phantom, unanimous-but-wrong,
clutter absorption, render ambiguity.}
\label{fig:taxonomy}
\end{figure}

\subsection{Knowledge-Grounded Query Evaluation}
\label{sec:exp-query}

We evaluate four conditions over identically constructed per-scene query
sets---48 items on the showroom, 32 on the robot hall (symbolic, explicit,
scene-level, and trap questions): (1) \emph{LLM-only} (Stage-A dump, no
verification), (2) \emph{ungated KG}, (3) \emph{verified-only KG}
(unverified nodes removed), and (4) \emph{gated KG} (unverified retained but
structurally demoted, Sec.~\ref{sec:query}). Each scene carries four trap
questions about objects that do not exist as labeled. On the showroom all
four traps target \emph{unverified-label} phantoms (keyboard, ladder); on
the robot hall two target an unverified-label phantom (refrigerator) and
two an \emph{escalation-verified} phantom (keyboard)---the latter probing
the gate's verifier-quality bound. A correct system must decline to assert
the trap labels. Note that LLM-only can ``block'' traps trivially: with no
map it tends to deny that specific objects exist at all (4/4 on the robot
hall, 2/4 on the showroom), so trap blocking is diagnostic only for the
grounded conditions.

\begin{table}[t]
\centering
\caption{Query benchmark, four conditions (accuracy / hallucination rate /
traps blocked), with abstention counts. Question sets: showroom $n{=}48$
(14 symbolic, 12 explicit, 9 implicit, 9 scene-level, 4 trap); robot hall
$n{=}32$ (15 symbolic, 12 explicit, 1 scene-level, 4 trap). An abstention
counts against accuracy (the denominator is always $n$) but is not a
hallucination; the verified-only condition's low accuracy is largely its
abstention cost.}
\label{tab:query}
\footnotesize
\setlength{\tabcolsep}{3pt}
\begin{tabular}{@{}lcccc@{}}
\toprule
 & LLM-only & Ungated & Verified-only & Gated \\
\midrule
Showroom ($n{=}48$) & 52.1\% / 20.8\% / 2/4 & 81.2\% / 18.8\% / 0/4 &
58.3\% / 16.7\% / 4/4 & \textbf{83.3\% / 16.7\% / 3/4} \\
\quad abstentions & 13 & 0 & 12 & 0 \\
Robot hall ($n{=}32$) & 50.0\% / 37.5\% / 4/4 & 56.2\% / 43.8\% / 0/4 &
46.9\% / 40.6\% / 2/4 & \textbf{59.4\% / 40.6\% / 2/4} \\
\quad abstentions & 4 & 0 & 4 & 0 \\
\bottomrule
\end{tabular}
\end{table}

On the showroom (Table~\ref{tab:query}, Fig.~\ref{fig:query}), grounding
alone (ungated) lifts
accuracy from 52.1\% to 81.2\% but blocks no traps; removing unverified
content (verified-only) blocks all four traps at a steep recall price (12
abstentions, 58.3\%); the gated condition is the balance point---83.3\% with
3/4 traps blocked. The two leaks are instructive and mechanistically
distinct. On the showroom, one \emph{count}-phrased trap slipped through
the quarantine itself: the model tallied a demoted node by reading its
\texttt{failedCandidateType} field, even though it correctly denied the
same object in yes/no phrasing. On the robot hall, the keyboard phantom is
\emph{escalation-verified}, so it passes any status-based gate---gating is
bounded by verifier quality; we report both as properties, not anomalies.
The robot hall reproduces the ordering at lower absolute levels because its
end-to-end label errors (banner$\rightarrow$refrigerator,
robot$\rightarrow$clutter) propagate into every grounded condition---an
honest end-to-end number rather than a curated one. We note the statistical
scope explicitly: at $n{=}48$/$32$ questions the gated-vs-ungated
\emph{accuracy} gap is a single question on each scene, so the benchmark is
diagnostic rather than powered for significance. The case for gating rests
on the mechanisms the benchmark isolates---trap blocking, the abstention
structure, and the direction of the hallucination change---and on the
downstream mission-level consequences of Sec.~\ref{sec:semnav}, not on the
accuracy deltas alone.

\begin{figure}[t]
\centering
\includegraphics[width=0.9\linewidth]{fig7_query_benchmark.png}
\caption{Four-condition query benchmark with trap-question panel.}
\label{fig:query}
\end{figure}

\subsection{Incremental Update Evaluation}
\label{sec:exp-update}

\paragraph{Synthetic scenarios.}
Controlled edits of the room-scoped robot hall isolate matcher behavior
(Fig.~\ref{fig:update}): \emph{(A) unchanged repeat}: 55/55 nodes carried,
zero VLM calls, zero identity switches; \emph{(B) move}: 1/1 moved with ID
preserved; \emph{(C) remove+insert}: exact absent/insert classification;
\emph{(D) occlusion}: at 20\% crop both affected objects still match (as
updated, labels kept), at 40\% matching fails into 2 false-absents + 2
false-news---the matcher's quantified breaking point, and the motivation for
overlap-based matching features in future work. On the combined-edit
re-scan (move, insertion, and removal applied together), selective
re-verification re-ran Stage B on only the two changed clusters: 26 of 822
VLM calls, \$0.23 versus \$7.37 for a full re-verification---3.1\% of
rebuild cost, and 0\% for scenario A. We scope this claim precisely: the
3.1\% figure comes from this controlled same-quality re-scan, where
unchanged clusters reproduce byte-identically. On the real
T2$\rightarrow$T3 transition the scan protocol itself changed
(higher-coverage acquisition), no cluster hash matched, and the epoch paid
full verification cost (\$3.54 for 41 clusters)---zero-cost carry-over is
a property of reproducible acquisition, not of the update mechanism alone.

\paragraph{Real three-epoch study.}
The three epochs of Table~\ref{tab:scenes} form a real update sequence:
T1 (manual, June) $\rightarrow$ T2 (exploration, early July) $\rightarrow$
T3 (99.6\%-coverage exploration with a controlled-change protocol, late
July). Cross-epoch registration is fully automatic
(FPFH~\cite{rusu2009fpfh}+RANSAC then ICP~\cite{besl1992icp};
5.1\,cm RMSE for T1$\rightarrow$T2, 4.2\,cm for T3$\rightarrow$T2). The
T1$\rightarrow$T2 diff records the room's exhibition-to-test-hall
reorganization (updated 4 / moved 3 / inserted 48 / absent 31; among the
moves, a fire extinguisher displaced 3.0\,m). This epoch also motivated the
generic-type exclusion in the move pass: without it, clutter piles matched
across 2--5\,m as fictitious ``moves'' (9 candidates, 3 credible).
The T2$\rightarrow$T3 diff (matched 9 / moved 1 / inserted 31 / absent 45,
Fig.~\ref{fig:epochs}) sits on top of measured natural drift: only 65\% of
in-room object geometry is mutually supported at 10\,cm between the two
scans two weeks apart---roughly a third of the room changed before any
protocol was applied.

\emph{Matching scorecard.} To score the real transitions without circular
reference to the matcher itself, we built an independent correspondence
reference by bidirectional point overlap between epochs (5\,cm
association, both directions $\ge$0.5 for a strong pair, 0.3--0.5
partial; the same construction as the T1 GT transfer of
Sec.~\ref{sec:exp-verify}), and replayed the recorded upserts against it.
The reference finds 4 strong and 6 partial static pairs for
T1$\rightarrow$T2 and 7 strong and 5 partial for T2$\rightarrow$T3.
Matching accuracy on strong pairs is 2/4 and 4/7; \textbf{identity
switches are zero on both transitions}---every miss is a
\emph{fragmentation} error in which re-segmentation changed an extent
beyond the 30\% tolerance, so the object was marked absent and
re-inserted under a new ID, never associated to the wrong object. The
per-miss diagnosis is uniform: all five failed the per-axis extent gate
(the most consequential case is the T3 ``new'' fire extinguisher, whose
cluster overlaps a T2 \texttt{clutter} cluster at 0.96--0.99 mutual
overlap but with width 0.11 vs.\ 0.28\,m). Strict false-new and
false-absent counts are 2/2 (T1$\rightarrow$T2) and 3/3
(T2$\rightarrow$T3), rising to 5/5 and 3/3 if partial pairs are also
demanded; nine split/merge cases (one cluster absorbing several) are
reported separately rather than forced into either bin. Multi-candidate
conflicts---the greedy resolution of Sec.~\ref{sec:matching}---occurred
zero times in pass~1 on both transitions and once in pass~2, where
nearest-centroid arbitration picked the owner-confirmed 0.64\,m chair
move. All three recorded \texttt{moved} claims are consistent with the
reference (a true move has no in-place overlap, so the reference can
refute but not confirm identity; the chair move is additionally
owner-confirmed). Repeating the reference construction at 10\,cm
association (twice the registration RMSE) leaves identity switches at
zero.

Against the controlled-change protocol, the system detected \textbf{2/2
physical changes at geometry level}: the moved chair is the epoch's
\emph{only} \texttt{moved} node (0.64\,m, ID preserved---a node matched
continuously since T1), and the introduced floor fan was correctly
\texttt{inserted} but typed \texttt{clutter} (error class (c) above), giving
1/2 at type level. Identity persistence is demonstrated across all three
epochs: the fire extinguisher, pump, industrial robot, and several chairs
retain their minted IDs through both reorganizations. The live graph has
since accumulated thirteen revisions through the same mechanism ---
rescans, owner refutations/restorations/corrections, re-split
decompositions, and finally a stale-node sweep in which eighteen
earlier-epoch detections with no support in the current scan were marked
absent; every transition, including the mistakes it corrects, is
reconstructible from node histories. The place layer
re-derives per-epoch names on the same SLIC regions
(Sec.~\ref{sec:place}): exhibition-era \texttt{machine\_access\_area} and
\texttt{operator\_control\_station} become clutter-dominated at T2 and
\texttt{machine\_operator\_station}/\allowbreak
\texttt{electrical\_services\_zone} at T3, at $\sim$\$0.03 per epoch.

\begin{figure}[t]
\centering
\includegraphics[width=0.9\linewidth]{fig8_incremental_update.png}
\caption{Controlled re-scan scenarios A--D (room-scoped robot hall, 55
objects), top-down view. Colored points are the segmented clusters, thin
rotated rectangles their explicit-model footprints with node names, and the
dark contour the test-room boundary; each edited object is annotated with
the matcher's decision. (A)~A repeat scan with no changes
carries all 55 verdicts at zero VLM cost. (B)~A 3.6\,m chair move is
classified \texttt{moved} with its node ID preserved. (C)~A removed object
is marked \texttt{absent} (node kept and flagged) and a newly placed object
is \texttt{inserted} under a new ID; only the insertion is re-verified. (D)~Partial occlusion: both cropped objects still match at 20\%,
but at 40\% matching fails into two false-absent plus two false-new
pairs---the matcher's quantified breaking point.}
\label{fig:update}
\end{figure}

\begin{figure}[t]
\centering
\includegraphics[width=0.9\linewidth]{fig9_two_epoch_update.png}
\caption{Real multi-epoch update, latest transition (T2$\rightarrow$T3),
top-down view. Colored points are the T3 segmented clusters inside the
test-room boundary; yaw-rotated node footprints are colored by diff state
(matched, moved with ID preserved, inserted, absent), with names on
specific-type objects. Oversized footprints of merged mega-clusters
(e.g., \texttt{machine\_012}, Sec.~\ref{sec:exp-ablation}) are faded.}
\label{fig:epochs}
\end{figure}

\subsection{Retrospective End-to-End Replay against Accumulated Owner
Ground Truth}
\label{sec:exp-endtoend}

After the corrective machinery of Secs.~\ref{sec:exp-seg}--\ref{sec:exp-update}
stabilized, we re-ran the entire pipeline from the raw scan with no owner
input at execution time and scored it against the live knowledge graph ---
by then a 46-object inventory embodying every owner restoration,
refutation, and correction. Because the corrective mechanisms were
developed from failures observed in this same environment, this
experiment is a \emph{retrospective end-to-end replay} on the
development scene, not an independent generalization test; what it
measures is how much of the owner-audited inventory the automated
layers recover once owner-informed development has concluded. The
replay runs the \emph{frozen} pipeline exactly as defined in
Secs.~\ref{sec:preprocess}--\ref{sec:corrective}---every trigger and
threshold fixed, no per-scene adjustment, no owner input at execution
time---so it is an independent \emph{run}, on a non-independent
\emph{scene}. The cafeteria (Sec.~\ref{sec:exp-ablation}) is the only
scene that is independent in both senses, and a multi-building
evaluation remains future work
(Sec.~\ref{sec:discussion}). One
corrective layer transfers mechanically and can therefore be checked
held-out: applying the hybrid structure filter to the cafeteria with
its T3 rules frozen (the occupancy-grid prior synthesized from the
scan's floor points) flags both ceiling-noise clusters with zero false
positives---both wall-flush doors survive, the trade-off that cost a
poster stand on T3---while the wall pillar evades the wall-hug test
and the seven carpet floor-noise clusters fall outside its rule
families entirely, consistent with the floor-material gap of
Sec.~\ref{sec:discussion}. The remaining corrective layers embed
scene-specific owner evidence by construction and are not claimed to
transfer. Table~\ref{tab:layers} reports cumulative geometric recall as
each automated layer is enabled. The verified path alone recovers 28\%;
diffuseness-gated re-splitting nearly doubles it; the wall-compound gate
(which self-triggers on the monitor wall and decomposes it into nine of
the ten owner-confirmed units) and the structure-condition ensemble bring
automated recall to 85\%. Type accuracy saturates far lower
($\sim$31\% of recalled objects) --- the label bottleneck is the
correlated verifier bias of Sec.~\ref{sec:exp-seg}, and closing it is
precisely the owner layer's role. For calibration within the lineage's
own reporting style, DK-SMF reports 78--81\% detection accuracy against
predefined ground-truth objects in its RGB-D environments~\cite{joo2025dksmf};
our 85\% is measured against an owner-audited inventory that includes
partially scanned and wall-flush objects, though the sensing modality
and ground-truth protocols differ. The re-run also reproduced the
deterministic backbone's stability end to end: the same diffuse webs,
the same compound decompositions, and a first-pass verification of the
fire extinguisher that had originally required owner restoration. The
precision column of Table~\ref{tab:layers} tracks what recall costs:
each added layer admits more candidate clusters, and the ensemble's
extra detection conditions temporarily dilute in-room precision (60\%
$\rightarrow$ 49\%) before the square-up step recovers part of it---
automated recall is bought with review burden, which is exactly the
owner-review efficiency that Sec.~\ref{sec:exp-owner} argues the model
should optimize.

\begin{table}[t]
\centering
\caption{Cumulative automated recall of the full re-run against the
owner-corrected 46-object graph (no owner input in the run itself).}
\label{tab:layers}
\small
\begin{tabular}{@{}lcc@{}}
\toprule
Automated layers (cumulative) & GT recall & In-room precision \\
\midrule
Verified path (Sec.~\ref{sec:decomposition}--\ref{sec:verification}) &
13/46 (28\%) & 37\% \\
+ diffuseness re-split & 26/46 (57\%) & 53\% \\
+ wall-compound gate & 36/46 (78\%) & 60\% \\
+ structure-condition ensemble & 37/46 (80\%) & 49\% \\
+ wall-compound panel square-up (Sec.~\ref{sec:corrective}, layer~2) & \textbf{39/46 (85\%)} & 51\% \\
\bottomrule
\end{tabular}
\end{table}

\subsection{Ablation and Sensitivity}
\label{sec:exp-ablation}

\paragraph{Escalation on/off.}
Disabling escalation reduces the pipeline to its base-4 fusion, already
tabulated as the base-4 column of Table~\ref{tab:verification}: escalation
contributes $+$11.1 points on the showroom and $+$6.3 on the robot hall,
at the bounded call cost analyzed in Fig.~\ref{fig:esccost}.

\paragraph{Gating style.}
Instructing the model to respect verification status blocks 1/4 traps;
structural demotion blocks 4/4 under identical prompts in this isolated
comparison (vs.\ 3/4 end-to-end, Sec.~\ref{sec:exp-query})---still the
single strongest argument that
reliability must live in the data contract, not in instructions.

\paragraph{Pose context.}
Injecting mounting-height context as a
targeted defense removed all three phantom ceiling keyboards in its pilot,
but applying it to every object \emph{reduced} escalated accuracy
(88.9\%$\rightarrow$75.0\%) and escalated precision
(100\%$\rightarrow$72.7\%)---a desk-height prior that helps ceilings hurts
everything else, so it is deployed as a targeted defense only.

\paragraph{Uni3D vocabulary.}
Swapping the Stage-A classifier vocabulary ranks generic-30 (66.7\%) above
curated-20 (48.7\%) and above the owner-derived 18-term set (33.3\%), with
$\sim$5-point run-to-run variance---narrowing the prompt set toward the
deployment vocabulary \emph{hurts} the embedding-level classifier,
confirming the domain gap that motivates VLM verification downstream.

\paragraph{Provisional-label ablation.}
Since the Stage-A label is supplied to the verifier as context, a fair
question is whether Stage B verifies or merely echoes it. We re-ran the
full showroom protocol (same renders, same verifier, base $+$ automatic
escalation) with the provisional label removed from both the system
prompt and the per-object context---dimensions and all perceptual
guardrails retained. The answer has two sides. With the label, Stage B
is a genuine but conservative corrector: on confirmed GT it departs from
Stage A on 8 of 36 objects, 5 correctly and 3 incorrectly, turning
Stage A's 77.8\% into 88.9\%. Without the label, the verifier is not an
independent classifier: decisiveness collapses (trigger rate
58\%\,$\rightarrow$\,90\%, unanimity 21\,$\rightarrow$\,5 objects) and
escalated accuracy falls to 38.9\%, with 22 of 26 departures from the
Stage-A label wrong. Two mechanisms account for the drop. Generic piles
lose their anchor: blind, the verifier names a specific type for
clutter (12 vs.\ 23 \texttt{clutter} answers against 22 GT), and on
specific-type objects accuracy still halves (10/14
$\rightarrow$ 4/14)---sparse renders under-determine the label, and the
prior selects among render-compatible hypotheses. The tier machinery,
however, degrades in the intended direction: blind uncertainty flows to
\texttt{unverified} (38 vs.\ 21 objects) rather than into confident
errors, though verified-tier precision does erode (unanimous 0.667,
escalated 0.5). The conclusion we draw is the one the architecture
already assumes: Stage~B is a \emph{verifier of hypotheses}, and the
provisional label is load-bearing context---the same anchoring power
that the candidate-shortlist experiment of Sec.~\ref{sec:exp-owner}
shows can mislead, which is why label adoption is agreement-gated and
owner-auditable rather than trusted outright.

\paragraph{Verifier model.}
Running the identical Stage-B protocol (Table~\ref{tab:vlmcmp}) across
five verifiers (two vendors, three tiers, two generations; reasoning
disabled for parity) makes escalation an economic instrument --- call
volume tracks verifier decisiveness monotonically (247 to 383 calls as
unanimous verifications fall from 26 to 9). Three findings. First,
within a generation the smaller tier splits more votes and resolves
fewer: its indecision inflates call volume by 14\% (383 vs.\ 335
calls), clawing back part of its nominal 3$\times$ per-token price
advantage --- and what survives is paid for in the lowest type
accuracy and the largest unverified residue (17 of 35 objects).
Second, all three current-generation
frontier models tie on type accuracy against the owner inventory
(7/13), but differ two-fold in the calls they need to get there ---
the same-generation mid-tier (Sonnet~5) is the most decisive (26
unanimous, 9 escalations), and the premium tier's extra deliberation
(18 escalations) buys no additional accuracy here. Third, the ranking
is generation-dominated, not vendor-dominated: both 2026-generation
mid-tiers beat both older Claude tiers, and cross-vendor differences
within the generation are smaller than the tier effects. Escalation
volume is thus a label-free, online proxy for verifier decisiveness on
the deployment's own data --- and a practical criterion for
deployment-time verifier selection under \emph{cost} constraints:
pick the cheapest verifier whose escalation rate has saturated. It is
not a proxy for accuracy: the upgrade check below shows the most
decisive model trading
unanimous precision for its decisiveness on the two full GT scenes, so
quality must still be priced by a small labeled probe set or by
accumulated owner feedback. The cost guide is robust where prices are
not: a vendor repricing days before our runs moved one verifier's
end-to-end cost five-fold (Table~\ref{tab:vlmcmp}, note~2), while the
decisiveness ordering it would be applied to is measured on the
deployment's own data.

\begin{table}[t]
\centering
\caption{Verifier-model comparison: identical Stage-B late fusion +
escalation on the same 35 in-room clusters, scored against the
owner-corrected graph.}
\label{tab:vlmcmp}
\small
\begin{tabular}{@{}lcccccc@{}}
\toprule
Verifier & Verified & Esc.-resolved & Unverified &
Type-correct$^{1}$ & Calls & Cost$^{2}$ \\
\midrule
Claude Sonnet 5 & \textbf{26} & 3 & \textbf{6} & \textbf{7/13} &
\textbf{247} & \$2.49 \\
GPT-5.6 (luna) & 21 & 7 & 7 & \textbf{7/13} & 287 & \$0.11 \\
Claude Opus 5 & 17 & 8 & 10 & \textbf{7/13} & 319 & \$5.28 \\
Claude Sonnet 4.6 & 15 & 6 & 14 & 6/13 & 335 & \$3.00 \\
Claude Haiku 4.5 & 9 & 9 & 17 & 5/13 & 383 & \$1.13 \\
\bottomrule
\end{tabular}

{\footnotesize $^{1}$ canonical-synonym type match on the clusters that
geometrically match a ground-truth object. $^{2}$ measured tokens
(in/out: 0.63M/40K, 0.42M/19K, 0.79M/53K, 0.74M/52K, 0.84M/59K, row
order) at vendor list rates of August 2026 (\$/Mtok in/out: 3/15,
0.20/1.20, 5/25, 3/15, 1/5). The GPT-5.6 row reflects the 2026-07-30
price reduction; at its launch rate (1/6) the same run costs \$0.54.
Sonnet~5's introductory rate (2/10) gives \$1.66.}
\end{table}

\paragraph{Verifier upgrade check.}
Because the ground truth is fixed and scoring is automatic, the headline
study of Sec.~\ref{sec:exp-verify} re-runs unchanged under a different
verifier. Re-running both scenes under Claude Sonnet~5---the most
decisive model above---confirms the decisiveness transfer: escalation
triggers fall from 58\% to 50\% (showroom) and from 51\% to 33\% (robot
hall), unanimous verifications rise from 21 to 25 and from 25 to 34, and
escalation calls drop from 232 to 200 and from 208 to 136. The freed
decisiveness is not free accuracy, however: unanimous precision falls
from 100\% to 91.7\% and from 94.4\% to 83.3\%, and escalated accuracy
lands at 69.4\% and 81.2\% (vs.\ 88.9\% and 84.4\%)---showroom deltas
several times the $\pm$3-point run-to-run variance measured in
Sec.~\ref{sec:exp-verify}, and directionally reproduced on the second
scene. The dominant new
failure is clutter absorption---real chairs, junction boxes, and a
control panel unanimously verified as \texttt{clutter}---i.e., the newer
model re-distributes error mass from indecision into confident
under-labeling on these 512-px render sets, while on the 1024-px renders
of Table~\ref{tab:vlmcmp} it ties for the best type accuracy.
Verifier upgrades therefore change \emph{where} errors live rather than
uniformly removing them, which is precisely why the framework treats
verification as tiered evidence behind a gate rather than as truth: the
tier structure, the escalation economics, and the owner loop survive a
verifier swap; the specific precision figures do not.

\begin{table}[!t]
\centering
\caption{Cafeteria stress scene: verification under verifier $\times$
render source (the identical 512-px multi-view protocol of
Sec.~\ref{sec:verification} throughout). Accuracy is escalated accuracy
on confirmed GT ($n{=}20$; the fragment-excluded column drops the five
confirmed sofa-fragment clusters, $n{=}15$); unanimous precision is over
confirmed unanimous votes. In every cell, 0 of the 7 sofa clusters are
recovered.}
\label{tab:cafe}
\small
\begin{tabular}{@{}llcccc@{}}
\toprule
Verifier & Render source & Esc.\ acc. & (excl.\ fragments) &
Unan.\ prec. & Trigger \\
\midrule
Sonnet 4-6 & 3\,cm voxel cloud & 35.0\% & 46.7\% & 57.1\% & 68\% \\
Sonnet 4-6 & full-resolution & 40.0\% & 53.3\% & 62.5\% & 60\% \\
Sonnet 5 & 3\,cm voxel cloud & 45.0\% & 60.0\% & 66.7\% & 40\% \\
Sonnet 5 & full-resolution & \textbf{60.0\%} & \textbf{80.0\%} &
\textbf{80.0\%} & 52\% \\
\bottomrule
\end{tabular}
\end{table}

\begin{figure}[!tp]
\centering
\includegraphics[width=\linewidth]{fig17_cafe_stress.png}
\caption{Cafeteria stress scene. (a)~The 28 clusters after automatic
yaw alignment and room crop; crosses mark the seven sofa
fragments/merges (absorbed as \texttt{clutter} in every condition of
Table~\ref{tab:cafe}), circles the seven carpet floor-noise clusters.
(b)~One cluster, rendered from the 3\,cm working cloud (left) and from
full-resolution re-associated points (right) under the identical
512-px protocol: the chair (top) is recovered by the render-source
change alone; the sofa fragment (bottom) is not. (c)~Escalated
accuracy (confirmed GT, $n{=}20$ per cell) and unanimous precision
(per-bar $n=8$--$15$ unanimous objects, annotated on each bar) across
verifier $\times$ render source; error bars are 95\% Wilson intervals,
wide on the precision bars because unanimous sets are small---the
point estimates order the cells, the intervals bound how far.}
\label{fig:cafe}
\end{figure}

\paragraph{Cross-floor stress scene.}
The cafeteria scene (Sec.~\ref{sec:setup}) probes the same protocol on
a furnishing domain the campaign never saw: carpeted floor, fabric sofa
booths, planter partitions (Fig.~\ref{fig:cafe}a). The backbone
transfers mechanically---57M points to 28 clusters in 13.5\,s after an
automatic wall-normal yaw alignment---but verification does not
(Table~\ref{tab:cafe}, first row): escalated accuracy reaches only
35\% against confirmed owner GT, and unanimous precision falls to
57.1\%, with two unanimous keyboard phantoms on carpet-noise clusters
and a chair unanimously absorbed as \texttt{clutter} (a fourth
unanimous miss, a table pedestal, falls on plausible-graded GT). The
errors are structured, not diffuse. The scene's
dominant object family---the sofas---enters the pipeline pre-broken,
split into five cushion fragments and two sofa--table merges, and every
one is absorbed (crosses in Fig.~\ref{fig:cafe}a); carpet pile leaves
seven residual floor-noise clusters above the RANSAC floor plane
(circles; the floor-material analogue of the ceiling remnants of
Sec.~\ref{sec:exp-seg}); and the confidence gate quarantines 9 of the
16 wrong labels over the 25 scored clusters as unverified while the
unanimous errors pass.

\paragraph{Render source $\times$ verifier.}
Two upgrades are then applied factorially (Table~\ref{tab:cafe},
Fig.~\ref{fig:cafe}c): re-rendering the identical 512-px view sets
from full-resolution re-associated points instead of the 3\,cm working
cloud (isolating render quality at unchanged protocol and cost,
$\sim$\$2.5/scene), and swapping the verifier for Sonnet~5. The two
levers act on different error classes. The render-source change
repairs \emph{whole-object} ambiguity for both verifiers: under
Sonnet~4-6 the keyboard phantoms vanish, both chairs recover
(including the unanimous miss; Fig.~\ref{fig:cafe}b, top), and both
doors resolve correctly, one via escalation---while the voxel renders
block the same chairs and doors under Sonnet~5 as well. The verifier
upgrade adds accuracy on top of whichever source it is given
(35$\rightarrow$45\% on voxel, 40$\rightarrow$60\% on full-resolution
sets): on this scene the newer model is uniformly better---in contrast
to its unanimous-precision loss on the industrial 512-px sets in the
upgrade check above---so even the \emph{sign} of a model upgrade is
domain-dependent. The render-source lever carries the same caveat: symmetric controls
that re-render \emph{both industrial scenes} from full-resolution
points reverse the effect---showroom escalated accuracy drops from
88.9\% to 80.6\% (base 77.8\%$\rightarrow$69.4\%; unanimous precision
holds at 100\%) and robot-hall from 84.4\% to 68.8\% (unanimous
precision 94.4\%$\rightarrow$83.3\%). The added texture invites
over-specific labels---a junction box read as a pipe assembly, the
wall TV as a ladder, clutter piles as keyboards, boxes, and trash
cans---where the coarser working-cloud render had left a correct
verdict, so the main campaign's voxel-source protocol stands as the
right default for the industrial scenes, and the cross-scene
comparison (first row of Table~\ref{tab:cafe} against
Table~\ref{tab:verification}) is same-protocol by construction. What
neither lever reaches is fragmentation: the seven sofa clusters are
absorbed in all four cells (Fig.~\ref{fig:cafe}b, bottom) and the
planters stay mislabeled.
Segmentation bounds what any verifier can see; render quality bounds
what a given verifier resolves; the model generation buys accuracy
only on top of both. Together with the upgrade check, the stress scene
fixes the reading of Table~\ref{tab:verification}: tier precision is a
scene- and domain-scoped measurement, to be re-estimated per
deployment---cheaply, via a small owner-confirmed probe set---rather
than a property of the architecture.

\paragraph{Segmentation parameters.}
The backbone's parameter sensitivity (36-combination sweep) and the
conditional refinement rules are ablated where their context lives, in
Sec.~\ref{sec:exp-seg}.

\paragraph{Voxel resolution and registration noise.}
Two controlled perturbations of the robot-hall backbone, scored per
object against the owner-grounded reference of
Sec.~\ref{sec:exp-seg}, bound the geometric operating regime. Voxel
size is not a free knob but part of a jointly calibrated triple with
$(\varepsilon, \mathrm{min\_points})$: at the baseline 3\,cm the full
backbone recovers 42/42 reference objects, but refining to 2\,cm
collapses recall to 0.36 (density rises, the fixed
$\varepsilon{=}0.30$ bridges gaps, 22 objects merge) and coarsening to
5\,cm collapses it to 0.33 (small objects starve below the 100-point
floor, 18 misses)---so the voxel size must be reported with the
clustering parameters, as Table~\ref{tab:sampleflow}'s pipeline
configuration does. Seeded per-point Gaussian noise applied to the raw
cloud before the backbone probes registration error at the point
level: at $\sigma{=}5$\,mm---the scale of the scanner's real per-scan
bundle error (4--5\,mm)---recall drops to 0.48, and by
$\sigma{=}20$\,mm (4$\times$ bundle error) the backbone effectively
fails (0.07), the dominant mechanism being saturation of the 5\,cm
RANSAC inlier band: un-removed wall residue grows from 267K to 853K
points and chains objects into merged clusters (39/42 merged) rather
than fragmenting them. We note the regime carefully: per-point scatter
at these levels is an artificial stress---the real $\sim$5\,cm
inter-epoch ICP residual is a rigid-transform error the pipeline
absorbs at the KG matching level (Sec.~\ref{sec:exp-update}), not
per-point noise---so the practical claim is robustness at the real
bundle-error scale, contingent on the structure-removal stage, which
is the first component to saturate. Color variation is not isolated by
any experiment in this paper: the render-source ablation above changes
point density and silhouette detail together with color, so it is an
appearance-quality control rather than a pure photometric one, and a
color-only perturbation study remains open.
